% GENERATED FILE. Edit canonical lessons, then run npm run generate:books.
% canonical-source-hash: fnv1a64:b68eef2633af8321
% canonical-lessons: UR-C200-pakora, UR-C200-qorma, UR-C200-achar, UR-C200-rayta, UR-C200-salad

\chapter{Pakora, Korma, Pickle, Raita, Salad}
\label{ch:ur-pakora-korma-pickle-raita-salad}
\hlchaptermodality{200}

\begin{chapteropening}
\textbf{By the end of this chapter:} \emph{I can say a pakora, korma, pickle, raita and a salad.}

Complete the last lesson of chapter 200: I can say a pakora, korma, pickle, raita and a salad.
\end{chapteropening}

\section[\texorpdfstring{\ur{پکوڑا} (pakoṛā)}{پکوڑا (pakoṛā)}]{\ur{پکوڑا} (pakoṛā) — a pakora}

\label{lesson:UR-C200-pakora}



\begin{quote}
Before the new one: say the Urdu for turmeric, then the Urdu for cumin.
\end{quote}

\subsection*{You'll want to know: \ur{پکوڑا}}
\textbf{\ur{پکوڑا}} — \emph{pakoṛā} — \textquotedblleft{}a pakora\textquotedblright{}.

\begin{grammarlens}[title={What you've built}]
One more word for this chapter.
\end{grammarlens}

\subsection*{Your turn}
\begin{itemize}
\raggedright
  \item \emph{Say it:} \emph{pakoṛā}
  \item \emph{Say it:} \emph{pakoṛā}, once more
  \item \emph{Say it:} say \emph{pakoṛā}
  \item \emph{Recall:} say the Urdu for turmeric, then the Urdu for cumin, then say \emph{pakoṛā} again
\end{itemize}

\begin{tcolorbox}[breakable,colback=teal!4,colframe=teal!35!black,title={Before you move on}]
What does \ur{پکوڑا} mean? (\textquotedblleft{}A pakora\textquotedblright{}.) Say it once more.
\end{tcolorbox}
\section[\texorpdfstring{\ur{قورمہ} (qormā)}{قورمہ (qormā)}]{\ur{قورمہ} (qormā) — korma}

\label{lesson:UR-C200-qorma}



\begin{quote}
Before the new one: say the Urdu for cumin, then the Urdu for a pakora.
\end{quote}

\subsection*{You'll want to know: \ur{قورمہ}}
\textbf{\ur{قورمہ}} — \emph{qormā} — \textquotedblleft{}korma\textquotedblright{}.

\begin{grammarlens}[title={What you've built}]
One more word for this chapter.
\end{grammarlens}

\subsection*{Your turn}
\begin{itemize}
\raggedright
  \item \emph{Say it:} \emph{qormā}
  \item \emph{Say it:} \emph{qormā}, once more
  \item \emph{Say it:} say \emph{qormā}
  \item \emph{Recall:} say the Urdu for cumin, then the Urdu for a pakora, then say \emph{qormā} again
\end{itemize}

\begin{tcolorbox}[breakable,colback=teal!4,colframe=teal!35!black,title={Before you move on}]
What does \ur{قورمہ} mean? (\textquotedblleft{}Korma\textquotedblright{}.) Say it once more.
\end{tcolorbox}
\section[\texorpdfstring{\ur{اچار} (achār)}{اچار (achār)}]{\ur{اچار} (achār) — pickle}

\label{lesson:UR-C200-achar}



\begin{quote}
Before the new one: say the Urdu for a pakora, then the Urdu for korma.
\end{quote}

\subsection*{You'll want to know: \ur{اچار}}
\textbf{\ur{اچار}} — \emph{achār} — \textquotedblleft{}pickle\textquotedblright{}.

\begin{grammarlens}[title={What you've built}]
One more word for this chapter.
\end{grammarlens}

\subsection*{Your turn}
\begin{itemize}
\raggedright
  \item \emph{Say it:} \emph{achār}
  \item \emph{Say it:} \emph{achār}, once more
  \item \emph{Say it:} say \emph{achār}
  \item \emph{Recall:} say the Urdu for a pakora, then the Urdu for korma, then say \emph{achār} again
\end{itemize}

\begin{tcolorbox}[breakable,colback=teal!4,colframe=teal!35!black,title={Before you move on}]
What does \ur{اچار} mean? (\textquotedblleft{}Pickle\textquotedblright{}.) Say it once more.
\end{tcolorbox}
\section[\texorpdfstring{\ur{رائتہ} (rāytā)}{رائتہ (rāytā)}]{\ur{رائتہ} (rāytā) — raita}

\label{lesson:UR-C200-rayta}



\begin{quote}
Before the new one: say the Urdu for korma, then the Urdu for pickle.
\end{quote}

\subsection*{You'll want to know: \ur{رائتہ}}
\textbf{\ur{رائتہ}} — \emph{rāytā} — \textquotedblleft{}raita\textquotedblright{}.

\begin{grammarlens}[title={What you've built}]
One more word for this chapter.
\end{grammarlens}

\subsection*{Your turn}
\begin{itemize}
\raggedright
  \item \emph{Say it:} \emph{rāytā}
  \item \emph{Say it:} \emph{rāytā}, once more
  \item \emph{Say it:} say \emph{rāytā}
  \item \emph{Recall:} say the Urdu for korma, then the Urdu for pickle, then say \emph{rāytā} again
\end{itemize}

\begin{tcolorbox}[breakable,colback=teal!4,colframe=teal!35!black,title={Before you move on}]
What does \ur{رائتہ} mean? (\textquotedblleft{}Raita\textquotedblright{}.) Say it once more.
\end{tcolorbox}
\section[\texorpdfstring{\ur{سلاد} (salād)}{سلاد (salād)}]{\ur{سلاد} (salād) — a salad}

\label{lesson:UR-C200-salad}



\begin{quote}
Before the new one: say the Urdu for pickle, then the Urdu for raita.
\end{quote}

\subsection*{You'll want to know: \ur{سلاد}}
\textbf{\ur{سلاد}} — \emph{salād} — \textquotedblleft{}a salad\textquotedblright{}.

\begin{grammarlens}[title={What you've built}]
One more word for this chapter.
\end{grammarlens}

\subsection*{Your turn}
\begin{itemize}
\raggedright
  \item \emph{Say it:} \emph{salād}
  \item \emph{Say it:} \emph{salād}, once more
  \item \emph{Say it:} say \emph{salād}
  \item \emph{Recall:} say the Urdu for pickle, then the Urdu for raita, then say \emph{salād} again
\end{itemize}

\begin{tcolorbox}[breakable,colback=teal!4,colframe=teal!35!black,title={Before you move on}]
What does \ur{سلاد} mean? (\textquotedblleft{}A salad\textquotedblright{}.) Say it once more.
\end{tcolorbox}
